\chapter{Complete Relations: D-pair, R-pair, PACK}

	\noindent
	The first think we are going to do is NOT ALWAYS using functions. We are going to use something that we will call ``Complete Relations''
	\\\\
	
	\noindent
	\textless ** \textit{See unnecesary comment 0004, to know the sad history about WHY they have that name} \textgreater
	\\\\
	
	\noindent
	Our relations could have more than one image element, for each element of the Domain. AND, they can have repetitions. They can have the same element of the image set assigned to diffrent elements of the Domain. The golden rule for this relations is that: \textbf{EVERY ELEMENT OF THE DOMAIN MUST HAVE, AT LEAST, ONE IMAGE}: no matter how... even cheating, repeating the use of elements of the Image set of the relation. They are called COMPLETE RELATIONS, because they always try to cover the Domain ``completely''.
	\\\\
	
	\noindent
	\textbf{R-pair}: means a pair of the relation. A t-upla of two elements, were left element belongs to the Domain, and the right element belongs to the Image.
	\\\\
	
	\noindent
	\textbf{D-pair}: two elements of the Domain, inside a t-upla of two elements. D-pairs belongs to the cartessian product ( Domain X Domain ).
	\\\\
	
	\noindent
	Imagine that we have two sets:\\
	$ A=\{a,b,c\}$\\
	$ B=\{1,2,3,4,5,6\}$\\
	And that we want to assign in our relation, 2,4,6 to a.\\
	Instead of creating the pair:\\
	$(a, \{2,4,6\})$\\
	We will add to the relation one R-pair, per each image:\\
	$(a,2)$\\
	$(a,4)$\\
	$(a,6)$
	\\\\
	
	\noindent
	This is a useless intermediate concept: The image set of an element of the domain. This set is made by all elements that are related to THAT element of the Domain, in some R-pair of the relation.The image set of a is $\{2,4,6\}$.
	\\\\
	
	\noindent
	\textbf{THE IMPORTANT ONE: A PACK}. A PACK of an element of the Domain, can be any possible subset, NOT EMPTY, proper or improper of the image set of that element. Examples for a:\\
	$ PACK(a)$ in r $= \{2,4,6\}$. VALID. \\
	$ PACK(a)$ in r $= \{2\}$. VALID.     \\
	$ PACK(a)$ in r $= \{6,2\}$. VALID.   \\
	$ PACK(a)$ in r $= \emptyset$. INVALID \\
	$ PACK(a)$ in r $= \{1,2\}$. INVALID
	\\\\
 	

\newpage
	
	
	